\documentclass{article}
\usepackage{mathtools} 
\usepackage{fontspec}
\usepackage[UTF8]{ctex}
\usepackage{amsthm}
\usepackage{mdframed}
\usepackage{xcolor}
\usepackage{amssymb}
\usepackage{amsmath}


% 定义新的带灰色背景的说明环境 zremark
\newmdtheoremenv[
  backgroundcolor=gray!10,
  % 边框与背景一致，边框线会消失
  linecolor=gray!10
]{zremark}{说明}

% 通用矩阵命令: \flexmatrix{矩阵名}{元素符号}{行数}{列数}
\newcommand{\flexmatrix}[4]{
  \[
  #1 = \begin{pmatrix}
    #2_{11}     & #2_{12}     & \cdots & #2_{1#4}   \\
    #2_{21}     & #2_{22}     & \cdots & #2_{2#4}   \\
    \vdots      & \vdots      & \ddots & \vdots     \\
    #2_{#31}    & #2_{#32}    & \cdots & #2_{#3#4}
  \end{pmatrix}
  \]
}

% 简化版命令（默认矩阵名为A，元素符号为a）: \quickmatrix{行数}{列数}
\newcommand{\quickmatrix}[2]{\flexmatrix{A}{a}{#1}{#2}}

\begin{document}
\title{习题 12.3}
\author{张志聪}
\maketitle

\section*{1}

\begin{itemize}
  \item (3)
        \begin{align*}
          \sum \frac{(-1)^n}{n^{p + \frac{1}{n}}}
           & = \sum \frac{(-1)^n}{n^p \cdot n^{\frac{1}{n}}}             \\
           & = \sum (-1)^n \frac{1}{n^p} \cdot \frac{1}{n^{\frac{1}{n}}} \\
        \end{align*}

        \begin{itemize}
          \item (i) $p < 0$；
                \begin{align*}
                  \lim\limits_{n \to \infty} \frac{1}{n^p} \cdot \frac{1}{n^{\frac{1}{n}}}
                   & = +\infty
                \end{align*}
                故级数发散。

          \item (ii) $p = 0$；
                \begin{align*}
                  \lim\limits_{n \to \infty} \frac{1}{n^0} \cdot \frac{1}{n^{\frac{1}{n}}}
                   & = \lim\limits_{n \to \infty} 1 \cdot \frac{1}{n^{\frac{1}{n}}} \\
                   & = 1
                \end{align*}
                故级数发散。

          \item (iii) $1 < p$；
                \begin{align*}
                  \lim\limits_{n \to \infty} \frac{\frac{1}{n^p} \cdot n^{\frac{1}{n}}}{\frac{1}{n^p}}
                   & = \lim\limits_{n \to \infty} n^{\frac{1}{n}} \\
                   & = 1
                \end{align*}
                因为$\sum \frac{1}{n^p}$收敛，由比较原则，级数绝对收敛。

          \item (iv) $0 < p \leq 1$；
                \begin{align*}
                  \lim\limits_{n \to \infty} \frac{\frac{1}{n^p} \cdot n^{\frac{1}{n}}}{\frac{1}{n^p}}
                   & = \lim\limits_{n \to \infty} n^{\frac{1}{n}} \\
                   & = 1
                \end{align*}
                因为$\sum \frac{1}{n^p}$发散，由比较原则，级数不绝对收敛。

                考虑条件收敛，对于交错级数$\sum (-1) \frac{1}{n^p}$，由莱布尼茨判别法可知其收敛。
                又有
                \begin{align*}
                  (\sqrt[x]{x})'
                   & = (e^{\ln\, \sqrt[x]{x}})'                                     \\
                   & = (e^{\frac{1}{x} \ln\, x})'                                   \\
                   & = e^{\frac{1}{x} \ln\, x} \left[\frac{1 - \ln\, x}{x^2}\right] \\
                   & = \sqrt[x]{x} \left[\frac{1 - \ln\, x}{x^2}\right]
                \end{align*}
                所以，$x > e$时，函数$\sqrt[x]{x}$单减，由归结原理可知，$n > 3$时，$\sqrt[n]{n}$单调且有界，
                因为改变级数的有限项并不影响原有级数的敛散性，因此利用阿贝尔判别法，
                级数收敛。
        \end{itemize}

  \item (5)

        由$\sum \frac{1}{n}$（调和级数）发散可知，该级数不绝对收敛。

        我们有
        \begin{align*}
          sum \frac{(-1)^n}{\sqrt{n}}
        \end{align*}
        收敛。
        \begin{align*}
          sum \frac{1}{n}
        \end{align*}
        发散。
        由极限的四则运算法则可知
        \begin{align*}
          sum \frac{(-1)^n}{\sqrt{n}} + \frac{1}{n}
        \end{align*}
        发散。

        注：这种拆分是有条件的，即知道拆开后的级数收敛或发散，本质上这是
        极限四则运算法则的应用，即：
        \begin{align*}
          \lim\limits_{n \to \infty} S_n
           & = \lim\limits_{n \to \infty} A_n + B_n                            \\
           & = \lim\limits_{n \to \infty} A_n + \lim\limits_{n \to \infty} B_n
        \end{align*}

  \item (8)
        \begin{align*}
          \sum \left|n! \left(\frac{x}{n}\right)^n\right|
          = \sum n! \left(\frac{|x|}{n}\right)^n
        \end{align*}
        于是
        \begin{align*}
          \lim\limits_{n \to \infty} \frac{|u_{n + 1}|}{|u_n|}
           & = \lim\limits_{n \to \infty} (n + 1)! \left(\frac{|x|}{n + 1}\right)^{n + 1} \cdot \frac{1}{n!} \left(\frac{n}{|x|}\right)^n \\
           & = \lim\limits_{n \to \infty} \left(\frac{n}{n+1}\right)^n \cdot |x|                                                          \\
           & = \lim\limits_{n \to \infty} \frac{|x|}{\left(\frac{n+1}{n}\right)^n}                                                        \\
           & = \lim\limits_{n \to \infty} \frac{|x|}{\left(1 + \frac{1}{n}\right)^n}                                                      \\
           & = \frac{1}{e} \cdot |x|
        \end{align*}

        (i) $|x|<e$，即$\frac{1}{e} \cdot |x| < 1$，由比式判别法，级数收敛，即绝对收敛。

        (ii) $|x| > e$，即$\frac{1}{e} \cdot |x| > 1$，由比式判别法可知，不绝对收敛。
        由于
        \begin{align*}
          c = \lim\limits_{n \to \infty} \frac{|u_{n+1}|}{|u_n|} > 1
        \end{align*}
        由保号性，存在$N_0$，$n > N_0$有
        \begin{align*}
          1 & < \frac{|u_{n+1}|}{|u_n|}
        \end{align*}
        于是可得
        \begin{align*}
          |u_{N_0}| < |u_n|  < |u_{n+1}|
        \end{align*}
        即当$n \to \infty$时，$u_n$的极限不可能为零，由柯西准则的推论可知，级数发散。

        (iii) $|x| = e$，即$\frac{1}{e} \cdot |x| = 1$。因为
        \begin{align*}
          \frac{|u_{n + 1}|}{|u_n|} & = \frac{|x|}{\left(1 + \frac{1}{n}\right)^n}
        \end{align*}
        由$\left(1 + \frac{1}{n}\right)^n$单增，可得$\frac{|x|}{\left(1 + \frac{1}{n}\right)^n}$单减，
        且收敛于1，
        于是对$N_0 > 1$，对所有的$n > N_0$有
        \begin{align*}
          \frac{|u_{n + 1}|}{|u_n|} \geq 1 \\
          |u_{n + 1}| \geq |u_n| \geq |u_{N_0}|
        \end{align*}
        即当$n \to \infty$时，$u_n$的极限不可能为零，由柯西准则的推论可知，级数发散。

  \item (9)

        \begin{align*}
          \sum |u_n| = \sum \frac{1}{n}
        \end{align*}
        发散，即不绝对收敛。

        令
        \begin{equation*}
          a_n =
          \begin{cases}
            1,  & n = 6k + 1, 6k + 2, 6k + 3 \\
            -1, & n = 6k + 4, 6k + 5, 6k + 6
          \end{cases}
          ,k = 0, 1, 2, ...
        \end{equation*}
        和
        \begin{align*}
          b_n = \frac{1}{n}
        \end{align*}
        所以$\{a_n\}$的部分和$A_n$有
        \begin{align*}
          |A_n| \leq 3
        \end{align*}
        且
        \begin{align*}
          \lim\limits_{n \to \infty} \frac{1}{n} = 0
        \end{align*}
        又$\{b_n\}$是单调减的，由A-D判别法，该级数条件收敛。

  \item (10)
        \begin{align*}
          \sum |u_n| = \sum \frac{1}{n}
        \end{align*}
        发散，即不绝对收敛。

        设级数的部分和为$S_n$，取它的子列$S_{3n}, \, n = 1, 2, 3, ...$，即
        \begin{align*}
          S_{3n}
           & = (1 + \frac{1}{2} - \frac{1}{3}) + \cdots + (\frac{1}{3n - 2} + \frac{1}{3n - 1} - \frac{2}{3n})
        \end{align*}
        令
        \begin{align*}
          T_n = 1 + \frac{1}{4} + \frac{1}{7} + \cdots + \frac{1}{3n - 2}
        \end{align*}
        因为
        \begin{align*}
          \frac{1}{3n - 1} - \frac{2}{3n}
        \end{align*}
        所以
        \begin{align*}
          T_n < S_{3n}
        \end{align*}
        $T_n$是$\sum \frac{1}{3n - 2}$的部分和数列，且该级数发散（与$\sum \frac{1}{n}$比较）。
        可知$T_n$发散且单增，所以
        \begin{align*}
          \lim\limits_{n \to \infty} T_n = +\infty
        \end{align*}
        于是
        \begin{align}
          \lim\limits_{n \to \infty} \{S_{3n}\} = +\infty
        \end{align}
        即$S_{3n}$发散。
        以为$S_{3n}$是$\{S_n\}$子列，所以$\{S_n\}$发散，即级数发散。

\end{itemize}

\section*{2}

\begin{itemize}
  \item (1)

        $\{a_n\}, a_n = \frac{x^n}{1 + x^n}$，$\sum b_n = \sum (-1)^n \frac{1}{n}$。

        $0 < x < 1$时，
        \begin{align*}
          \lim\limits_{n \to \infty} a_n
           & = \lim\limits_{n \to \infty} \frac{x^n}{1 + x^n}         \\
           & = \lim\limits_{n \to \infty} \frac{1}{\frac{1}{x^n} + 1} \\
           & = 0
        \end{align*}
        且$\{a_n\}$单减，即单调有界。

        $x = 1$时，
        \begin{align*}
          \lim\limits_{n \to \infty} a_n
           & = \lim\limits_{n \to \infty} \frac{1}{1 + 1} \\
           & = \frac{1}{2}
        \end{align*}
        单调有界。

        $x > 1$时，且
        \begin{align*}
          \lim\limits_{n \to \infty} a_n
           & = \lim\limits_{n \to \infty} \frac{x^n}{1 + x^n}         \\
           & = \lim\limits_{n \to \infty} \frac{1}{\frac{1}{x^n} + 1} \\
           & = 1
        \end{align*}
        $\{a_n\}$单增，即单调有界。

        又因为$\sum b_n$是收敛的，由A-D判别法，该级数收敛。

  \item (2)

        利用积化和差公式
        \begin{align*}
          2sin\, x \, sin \, y = cos(x - y) - cos(x + y)
        \end{align*}
        考虑
        \begin{align*}
          2 sin \frac{x}{2} \sum\limits_{k = 1}^n sin kx
           & = 2 sin \frac{x}{2} \left( sin x + sin 2x + \cdots sin nx \right)                        \\
           & = 2 sin \frac{x}{2} sin x + 2 sin \frac{x}{2} sin 2x + \cdots + 2 sin \frac{x}{2} sin nx \\
           & = cos \frac{x}{2} - cos \frac{3x}{2} + cos \frac{3x}{2} - cos \frac{5x}{2}
          + \cdots + \cos\frac{(n - \frac{1}{2})x}{2} - cos \frac{(n + \frac{1}{2})x}{2}              \\
           & = cos \frac{x}{2} - cos \frac{(n + \frac{1}{2})x}{2}
        \end{align*}
        于是
        \begin{align*}
          \sum\limits_{k = 1}^n sin kx = \frac{cos \frac{x}{2} - cos \frac{(n + \frac{1}{2})x}{2}}{2 sin \frac{x}{2}}
        \end{align*}
        故
        \begin{align*}
          \frac{cos \frac{x}{2} - 1}{2 sin \frac{x}{2}} \leq \sum\limits_{k = 1}^n sin kx \leq \frac{cos \frac{x}{2} + 1}{2 sin \frac{x}{2}}
        \end{align*}
        注：$\frac{x}{2} \in (0, \pi)$，$sin \frac{x}{2} \neq 0$。

        综上，$\sum sin nk$的部分和数列有界。

        另外
        \begin{align*}
          \lim\limits_{n \to \infty} \frac{1}{n^\alpha} = 0
        \end{align*}
        且单减。

        故，利用A-D判别法，级数收敛。

  \item (3)
        \begin{align*}
          \sum (-1)^n \frac{cos^2\, n}{n}
           & = \sum (-1)^n \frac{1 + cos 2n}{2n}                        \\
           & = \sum (-1)^n \frac{1}{2n} + \sum (-1)^n \frac{cos 2n}{2n}
        \end{align*}
        由莱布尼茨判别法可知$\sum (-1)^n \frac{1}{2n}$收敛。

        利用积化和差公式
        \begin{align*}
          2cos\, x \, cos\, y = cos (x - y) + cos (x + y)
        \end{align*}
        考虑
        \begin{align*}
          2 cos 1 \sum\limits_{k=1}^n (-1)^k cos 2k
           & = 2 cos 1 [-cos 2 + cos 4 - cos 6 + \cdots + (-1)^n cos (2n) ]                      \\
           & = - 2 cos 1 cos 2 + 2 cos 1 cos 4 - 2 cos 1 cos 6 + \cdots + 2cos 1 (-1)^n cos (2n) \\
           & = -(cos 1 + cos 3) + (cos 3 + cos 5) - (cos 5 + cos 7)                              \\
           & + \cdots + (-1)^n [cos (2n - 1) + cos (2n + 1)]                                     \\
           & = - cos 1 + (-1)^n cos (2n + 1)                                                     
        \end{align*}
        所以
        \begin{align*}
          \sum\limits_{k=1}^n (-1)^k cos 2k = \frac{- cos 1 + (-1)^n cos (2n + 1)}{2 cos 1}
        \end{align*}
        故
        \begin{align*}
          \frac{-1 - cos 1}{2 cos 1} \leq \sum\limits_{k=1}^n (-1)^k cos 2k \leq \frac{1 - cos 1}{2 cos 1}
        \end{align*}
        所以，$\sum\limits_{k=1}^n (-1)^k cos 2k$的部分和有界。

        又因为
        \begin{align*}
          \lim\limits_{n \to \infty} \frac{1}{2n} = 0
        \end{align*}
        且$\{\frac{1}{2n}\}$单减。
        利用A-D判别法，级数$\sum (-1)^n \frac{cos 2n}{2n}$收敛。

        综上
        \begin{align*}
          \sum (-1)^n \frac{cos^2\, n}{n} & = \sum (-1)^n \frac{1}{2n} + \sum (-1)^n \frac{cos 2n}{2n}
        \end{align*}
        收敛。

\end{itemize}


\end{document}